% PHỤ LỤC: A. Phân công nhiệm vụ; B. Mã nguồn và tái lập.
\clearpage
\phantomsection
\section*{\hypertarget{ap:A}{}PHỤ LỤC A. PHÂN CÔNG NHIỆM VỤ}
\addcontentsline{toc}{section}{Phụ lục A. Phân công nhiệm vụ}

\begin{table}[H]
\centering\small
\begin{tabular}{|p{3.6cm}|p{10.6cm}|p{1.9cm}|}
\toprule
Thành viên & Công việc & Đóng góp \\
\midrule
Lý Anh Quân & Tổ chức mã nguồn Python và kiểm thử tự động & 20\% \\
Vũ Trần Hoàng Quân & Kiểm chứng và viết Chương 4: chuỗi dừng và mô hình ARMA & 20\% \\
Hồ Lê Ngọc Quang & Kiểm chứng và viết Chương 5: chuỗi không dừng và ARIMA & 20\% \\
Phạm Nhật Quang & Kiểm chứng và viết Chương 6: dự báo và đánh giá dự báo & 20\% \\
Phạm Anh Quốc & Soạn Chương 1 đến 3, thảo luận, kết luận và dàn trang LaTeX & 20\% \\
\bottomrule
\end{tabular}
\end{table}

\clearpage
\phantomsection
\section*{\hypertarget{ap:B}{}PHỤ LỤC B. MÃ NGUỒN VÀ TÁI LẬP}
\addcontentsline{toc}{section}{Phụ lục B. Mã nguồn và tái lập}
\setcounter{table}{0}\renewcommand{\thetable}{B.\arabic{table}}


\subsection*{B.1. Notebook đính kèm}

Tệp \texttt{AS3179\_Nhom7.ipynb} chứa mã xử lý dữ liệu, các hàm thống kê, mô phỏng, ước lượng, dự báo và mã sinh bảng, hình của báo cáo. Trong bản PDF, có thể lưu notebook bằng biểu tượng kẹp giấy dưới đây hoặc từ danh sách tệp đính kèm của trình đọc PDF. Trên LMS, cùng tệp này được nộp riêng để có thể mở trực tiếp bằng Jupyter.

\par\noindent
\immediate\pdfobj stream attr {/Type /EmbeddedFile } file {\notebookdir/AS3179_Nhom7.ipynb}
\edef\notebookstream{\the\pdflastobj}
\immediate\pdfobj{<< /Type /Filespec /F (AS3179_Nhom7.ipynb) /UF (AS3179_Nhom7.ipynb)
 /Desc (AS3179 - Nhom 7 - Jupyter Notebook) /EF << /F \notebookstream\space 0 R >> >>}
\edef\notebookfile{\the\pdflastobj}
\immediate\pdfobj{<< /Names [(AS3179_Nhom7.ipynb) \notebookfile\space 0 R] >>}
\edef\notebooknames{\the\pdflastobj}
\pdfnames{/EmbeddedFiles \notebooknames\space 0 R}
\leavevmode\pdfannot width 12pt height 12pt depth 0pt {/Subtype /FileAttachment /FS \notebookfile\space 0 R /Name /Paperclip /Contents (AS3179_Nhom7.ipynb) /C [0.1 0.3 0.6]}
\hspace{16pt}\texttt{AS3179\_Nhom7.ipynb}\quad (mã nguồn của Nhóm 7)

\subsection*{B.2. Chạy từ năm tệp dữ liệu gốc}

\begin{enumerate}
\item Đặt notebook cạnh \texttt{AirPassengers.csv}, \texttt{EuStockMarkets.csv}, \texttt{LakeHuron.csv}, \texttt{lynx.csv} và \texttt{Nile.csv}. Cũng có thể giữ nguyên cấu trúc \texttt{data/raw/} trong kho dự án. Notebook kiểm tra SHA-256 của từng CSV trước khi dùng.
\item Chọn môi trường Python 3.12 trong Jupyter. Có thể cài các thư viện của lần chạy báo cáo bằng lệnh sau trong môi trường đó.
\begin{codeblock}
python -m pip install jupyter ipykernel pillow numpy==2.5.5 pandas==3.1.1 scipy==1.18.1
python -m pip install matplotlib==3.11.2 statsmodels==0.15.0
\end{codeblock}
\item Mở notebook, chọn \textit{Restart Kernel and Run All Cells}. Đợi chạy xong cả phần Bài 12. Các ô đầu tiên nạp trực tiếp mã của nhóm vào bộ nhớ; các ô tiếp theo đọc dữ liệu, thực hiện kiểm chứng và lưu kết quả. Hạt giống mô phỏng là \texttt{3179}, với dòng số ngẫu nhiên riêng cho từng thí nghiệm.
\end{enumerate}

\subsection*{B.3. Tìm phép tính và đối chiếu đầu ra}

\noindent\begin{tabularx}{\linewidth}{@{}p{4.0cm}X@{}}
\toprule
Phần trong notebook & Nội dung trong báo cáo \\
\midrule
Dữ liệu và khám phá & Sáu chuỗi HU, LY, DR, NI, AP, DP; bảng mô tả và đồ thị ở Chương 3. \\
Bài 11, mục 11.1--11.14 & Mười bốn nội dung ARMA của Chương 4. Bảng cuối Bài 11 ghi địa chỉ báo cáo cho từng mục. \\
Bài 12, mục 12.1--12.14 & Mười bốn nội dung ARIMA và dự báo của Chương 5 và 6. Chẳng hạn, mục 12.13 tính kiểm định Diebold--Mariano ở mục 6.3(c) cho AP, NI và DP. \\
\bottomrule
\end{tabularx}

Sau một lần chạy, \texttt{data/processed/} có sáu CSV, \texttt{outputs/tables/} có 67 bảng (mỗi bảng ở ba dạng CSV, Markdown và LaTeX) và \texttt{outputs/figures/} có 26 hình PNG. Năm tệp gốc, sáu tệp đã xử lý, notebook, PDF và tệp LaTeX được nộp rời. Để dựng lại bản PDF, đặt \texttt{main.tex} cạnh notebook, giải nén \texttt{figures.zip} tại đó và chạy \texttt{pdflatex main.tex} cho đến khi các tham chiếu ổn định.
